\rhead{CS101: Intro to Programming}
\phantomsection 
\section*{Introduction to Programming (CS101)}
\label{course:CS101}

\noindent \small{\hyperref[main:scheme]{$\leftarrow$ Back to Scheme}} \vspace{0.5cm}

\noindent \textbf{Credits:} 4 (4-0-0)

\subsection*{Course Content}
\noindent\textbf{Unit 1} \\
\textit{Problem Solving and Core Concepts:} Introduction to Algorithms and Flowcharts, Structure of a C++ Program, Preprocessor directives, Console Input/Output (iostream), Fundamental Data Types and Literals, Variable Scope and Initialization, Type Conversions (Implicit/Explicit), Operators (Arithmetic, Relational, Logical), Control Flow (Conditionals, Loops, break/continue).

\vspace{0.3cm}

\noindent\textbf{Unit 2} \\
\textit{Compound Data Types and Memory:} Pointers (Declaration, Dereferencing), Dynamic Memory Allocation (new/delete), Array structures (1D/2D C-style Arrays vs std::vector), String manipulation (C-style vs std::string), User-Defined Types (Structures for grouping data, Unions for memory sharing).

\vspace{0.3cm}

\noindent\textbf{Unit 3} \\
\textit{Modular Programming and OOP Fundamentals:} Function definitions and Overloading, Parameter Passing mechanisms (Value, Pointer, Reference), Introduction to Classes and Objects, Encapsulation principles, Access Specifiers (public, private, protected), Object lifecycle management (Constructors, Destructors).

\vspace{0.3cm}

\noindent\textbf{Unit 4} \\
\textit{Advanced Object-Oriented Concepts:} Inheritance types (Base and Derived classes), Code reusability, Polymorphism (Compile-time Overloading vs Run-time Overriding), Virtual Functions, Abstract Classes, Static Data Members, Exception Handling (try, catch, throw), Standard Exception hierarchy.

\vspace{0.3cm}

\noindent\textbf{Unit 5} \\
\textit{STL and Modern C++:} Standard Template Library (Sequence Containers, Associative Containers, Container Adaptors), Iterators and Generic Algorithms (Sort, Find), Modern C++ features (Smart Pointers for memory safety, Lambda Expressions, Range-based loops), Utilities (std::pair, std::tuple).

\newpage
\rhead{Curriculum Schema}